%% ───────────────────────────────────────────────────────────────────
\chapter{Gli Obiettivi di un Sistema Pensionistico: Quadro Teorico}
%% ───────────────────────────────────────────────────────────────────

%% ─── PROMPT LLM ─────────────────────────────────────────────────
%% Scrivi il Capitolo 1 di una tesi magistrale in economia pubblica.
%% Obiettivo: fornire il framework teorico che guiderà tutta la tesi.
%% Tono: rigore scientifico + prosa scorrevole. Citare Bosi e Artoni
%% come fonti teoriche primarie, Rosen-Gayer-Rapallini (McGraw-Hill 2023)
%% come manuale di riferimento del corso.
%%
%% IMPORTANTE: questo capitolo deve stabilire il "trilemma previdenziale"
%% (sostenibilità / adeguatezza / equità) come framework analitico
%% che verrà usato in tutta la tesi per valutare riforme e strategie.
%% ────────────────────────────────────────────────────────────────

\section{Le tre funzioni del sistema pensionistico pubblico}

    %% PROMPT: Spiega le tre funzioni (previdenziale, assicurativa,
    %% assistenziale) con definizioni precise da manuale.
    %% Per ciascuna: definizione tecnica + esempio concreto + quale
    %% tipologia di pensione vi corrisponde (vecchiaia, invalidità,
    %% sociale, superstiti). Concludi mostrando che il settore privato
    %% può svolgere queste funzioni solo parzialmente e imperfettamente
    %% (risparmio individuale, famiglia, volontariato) -- giustificando
    %% così l'intervento pubblico.
    \subsection{Funzione previdenziale: garantire un tenore di vita
                simile a quello dell'ultima fase lavorativa}
    \subsection{Funzione assicurativa: protezione dal rischio
                di riduzione o fluttuazione del reddito nel tempo}
    \subsection{Funzione assistenziale: reddito minimo
                per un'esistenza dignitosa}
    \subsection{Il privato può sostituire il pubblico?
                Risparmio individuale, famiglia e volontariato
                come alternative imperfette}

\section{Le ragioni dell'intervento pubblico obbligatorio}

    %% PROMPT: Spiega le quattro giustificazioni dell'intervento
    %% pubblico OBBLIGATORIO e UNIVERSALE (non solo facoltativo).
    %% 1. Fallimenti del mercato assicurativo privato:
    %%    - asimmetrie informative sul rischio degli investimenti
    %%    - assenza di copertura dal rischio inflazione
    %%    - assenza di economie di scala (costi di marketing/gestione)
    %% 2. Miopia individuale (bene meritorio / paternalismo):
    %%    gli individui sottostimano il rischio di povertà in vecchiaia
    %%    e risparmiano meno del socialmente ottimale.
    %% 3. Azzardo morale: senza obbligo, un individuo potrebbe
    %%    non risparmiare sapendo che lo Stato interverrà comunque.
    %% 4. Selezione avversa: senza universalità, solo i lavoratori
    %%    meno produttivi rimarrebbero nel sistema, riducendo il monte
    %%    contributivo e portando alla crisi del sistema.
    %% Fonti: Bosi, Artoni, Rosen-Gayer-Rapallini cap. 12.
    \subsection{Fallimenti del mercato assicurativo privato:
                selezione avversa, asimmetrie informative
                e assenza di economie di scala}
    \subsection{Miopia individuale e bene meritorio:
                il paternalismo come giustificazione
                dell'obbligatorietà}
    \subsection{Azzardo morale e universalità:
                perché senza contribuzione obbligatoria
                il sistema si sgretola}
    \subsection{Il fallimento informativo come giustificazione
                aggiuntiva: perché la maggioranza dei lavoratori
                non conosce il proprio tasso di sostituzione atteso}

\section{I modelli di finanziamento: PAYG vs Capitalizzazione}

    %% PROMPT: Presenta i tre modelli con rigore tecnico.
    %% Per ciascun modello: meccanismo di funzionamento,
    %% vantaggi, svantaggi, rischi specifici.
    %%
    %% MODELLO 1 -- Capitalizzazione (FF: Fully Funded):
    %%   ottica assicurativa individuale, rendimento = r del mercato.
    %%   Pensione: P_FF = c * S_t * (1+r)
    %%   Rischio: volatilità dei mercati finanziari.
    %%
    %% MODELLO 2 -- Ripartizione (PAYG: Pay-As-You-Go):
    %%   patto intergenerazionale, rendimento implicito = (1+m)(1+n).
    %%   Pensione: P_PAYG = c * S_t * (1+m)(1+n)
    %%   Rischio: transizione demografica (riduzione di n).
    %%   I due sistemi danno la stessa pensione se (1+r) ≈ (1+m+n)
    %%   [Teorema di Aaron].
    %%
    %% MODELLO 3 -- NDC (Notional Defined Contribution):
    %%   ibrido: PAYG nella sostanza, DC nella forma.
    %%   Il conto è "virtuale" (nessuna riserva reale),
    %%   l'interesse implicito g = media mobile 5 anni del PIL nominale.
    %%   Spiega perché è una "terza via" e i suoi vantaggi
    %%   rispetto ai due modelli puri.
    %%
    %% Aggiungi: gli effetti economici della previdenza PAYG sul
    %% comportamento individuale (effetto sostituzione della ricchezza,
    %% effetto anticipo pensionamento, effetto eredità).
    \subsection{Il sistema a capitalizzazione (FF):
                ottica assicurativa individuale e rischio finanziario}
    \subsection{Il sistema a ripartizione (PAYG):
                il patto intergenerazionale e il teorema di Aaron}
    \subsection{Il Notional Defined Contribution (NDC):
                la terza via ibrida -- PAYG nella sostanza,
                DC nella forma}
    \subsection{Gli effetti economici della previdenza PAYG:
                spiazzamento del risparmio privato, effetto anticipo
                pensionamento ed effetto eredità}

\section{Il trilemma previdenziale: il framework della tesi}

    %% PROMPT: Presenta il "trilemma sostenibilità–adeguatezza–equità"
    %% come il prisma analitico che guiderà tutta la tesi.
    %% Definisci con precisione i tre vertici:
    %%
    %% 1. SOSTENIBILITÀ FINANZIARIA:
    %%    equazione macroeconomica estesa:
    %%    S_p/PIL = (N_p/N_a) × (N_a/POP) × (POP/N_l) × [(S_p/N_p)/(PIL/N_l)]
    %%    Spiega i quattro fattori: istituzionale, demografico,
    %%    occupazionale, redistributivo (tasso di sostituzione medio).
    %%
    %% 2. ADEGUATEZZA:
    %%    tasso di sostituzione individuale (pensione / ultimo salario),
    %%    soglia di povertà relativa in età anziana,
    %%    concetto di "tenore di vita dignitoso" (funzione previdenziale).
    %%
    %% 3. EQUITÀ INTERGENERAZIONALE:
    %%    tasso di rendimento implicito (IRR) per generazione,
    %%    confronto Baby Boomer vs Millennials/Gen Z,
    %%    concetto di neutralità attuariale.
    %%
    %% Concludi spiegando la TESI CENTRALE: le riforme Dini e Fornero
    %% hanno risolto il vertice della sostenibilità ma hanno creato
    %% un problema crescente sui vertici adeguatezza ed equità.
    \subsection{Sostenibilità finanziaria: la decomposizione
                dell'equazione $S_p/PIL$ nei quattro fattori}
    \subsection{Adeguatezza delle prestazioni: tasso di sostituzione,
                pensione minima e povertà relativa in età anziana}
    \subsection{Equità intergenerazionale: il tasso di rendimento
                implicito (IRR) e la neutralità attuariale}
    \subsection{Il trade-off al centro della tesi: sostenibilità
                comprata al prezzo dell'adeguatezza futura}
